\chapter{Literature Review}

\section{Brief}
This chapter examines the literature on how the design of online messaging platforms shapes users’ ability to understand, manage, and revoke consent, with a specific focus on photo sharing on WhatsApp. It begins with a brief overview of digital messaging and photo-based communication to situate the social and privacy-sensitive nature of image sharing. It then reviews relevant theoretical and consent-related perspectives to explain why consent revocation is not only a technical issue but also a contextual and interpersonal one. The chapter subsequently examines research on instant messaging, photo sharing, and interface design to assess how current platform features support or constrain meaningful user control. Finally, it synthesises the literature to identify research gaps that motivate the proposed user-centred revocation design.


\section{Overview of Digital Messaging and Photo-Based Communication}
Digital messaging has developed from earlier forms of computer-mediated communication into mobile, app-based, and increasingly multimodal communication environments, where text, images, and other media are integrated within the same platforms. Herring’s historical overview highlights this shift from predominantly text-based CMC toward convergent and multimodal communication systems, providing a useful background for understanding contemporary messaging platforms \cite{Herring2019Coevolution}. Research on mobile messaging also shows how texting became embedded in everyday social contact, particularly in peer communication shaped by social norms and situational constraints \cite{ito_mobile_2005}. In parallel, early camera-phone studies show that images were used not only for personal capture but also for communication with absent others and remote tasks, helping establish photo messaging as a routine social practice \cite{kindberg_i_2005}. This transition is important because it created communication settings in which image sharing carries social meaning and privacy expectations that may later complicate consent revocation.


\section{Digital Privacy in Modern Communication}
This section grounds the dissertation in theories that explain why consent revocation is not simply a delete button, but a socially negotiated and context-dependent action.

\subsection{Communication Privacy Management (CPM)}
Communication Privacy Management (CPM) theory explains privacy as an ongoing process of boundary management rather than a fixed state, making it useful for analysing disclosure in interpersonal and mediated communication \cite{child_petronio_cmc_privacy_2011, petronio_cpm_2020}. Petronio and Child (2020) describe CPM as a framework in which individuals treat private information as something they own and control, and they manage access to it through privacy rules shaped by context, relationship, perceived risk, and expectations about appropriate use after disclosure \cite{petronio_cpm_2020}. A key concept is co-ownership: once information is shared, recipients may become legitimate co-owners who develop their own assumptions about how that information can be retained, discussed, or circulated \cite{petronio_cpm_2020}. CPM therefore shifts attention from information itself to the negotiated rules and expectations that organise its flow. 

This matters for consent revocation in messaging because consent revocation in messaging is not only a technical question of whether content can be deleted; it is also a social question about whether previously established privacy expectations can be changed after disclosure \cite{child_petronio_cmc_privacy_2011, petronio_cpm_2020}. CPM helps explain why post-sharing control is often contentious: a sender may understand revocation as a reassertion of personal control, while a recipient may interpret the same act as changing the rules after access has already been granted \cite{petronio_cpm_2020}. Hence, CPM provides a strong theoretical foundation for studying consent and control in messaging-based photo sharing, where interpersonal expectations are central to how users judge fairness, trust, and appropriateness.

Applied to revocable or dynamic consent in messaging, CPM suggests that revocation should be understood as an attempt to re-coordinate privacy boundaries rather than simply erasing data \cite{petronio_cpm_2020}. The theory’s concept of boundary turbulence is especially relevant here, as turbulence occurs when co-owners hold conflicting privacy rules or when expected information flows break down \cite{petronio_cpm_2020}. In messaging environments such as WhatsApp-style photo sharing, this framework helps explain why consent may need to be updated over time and why platform mechanisms for revocation must operate within the tension between social boundary expectations and the technical ease of retention, copying, or redistribution \cite{child_petronio_cmc_privacy_2011, petronio_cpm_2020}.


\subsection{Contextual Integrity (CI)}
Contextual Integrity (CI) reframes privacy as appropriateness of information flows, not simply secrecy or access control \cite{nissenbaum_contextual_2011}. Nissenbaum (2011) argues that information sharing is governed by contextual norms: who is involved (actors), what is shared (attributes), and how it moves (transmission principles) \cite{nissenbaum_contextual_2011}. This is directly relevant to WhatsApp photo sharing because users often treat a private chat as a bounded context with an implied “right audience,” even when the content is not technically protected from recipient-side redistribution. Under CI, harm occurs when a photo’s flow changes context (e.g., from a one-to-one chat into a group, another platform, or offline circulation) because the transfer violates the original contextual expectations, even if the sender originally shared willingly \cite{nissenbaum_contextual_2011}.

CI is also useful for critiquing the limits of the notice-and-consent paradigm in everyday messaging. Nissenbaum (2011) critiques notice-and-consent as insufficient for online privacy governance, while Bourgeus and Vandercruysse (2024) argue that overreliance on consent often fails to produce meaningful privacy control in complex data ecosystems and motivate contextual and communal alternatives \cite{nissenbaum_contextual_2011,bourgeus_vandercruysse_2024}. This matters for photo sharing because images frequently involve group privacy: a photo can affect the sender, the recipient, and any depicted co-subjects, each with potentially different expectations about appropriate sharing \cite{bourgeus_vandercruysse_2024}. Design implication: revocation mechanisms should be framed as restoring context-appropriate flow (limiting further redistribution and future access within WhatsApp), and the UI should represent that permissions apply to a social context (audience and transmission rules), not just to an individual recipient or a one-time send action.

\subsection{Media Richness Theory (MRT) and Social Presence}

Media Richness Theory (MRT) holds that communication media vary in their capacity to convey information in ways that support interpretation and mutual understanding, and that richer media are generally better suited to tasks involving ambiguity (equivocality) \cite{daft_lengel_mrt_1986}. In Daft and Lengel’s account, rich media allow more immediate feedback and convey a greater variety of cues, which can help communicators clarify meaning and reduce misunderstanding when issues are socially sensitive or open to multiple interpretations \cite{daft_lengel_mrt_1986}. WhatsApp-style messaging that includes photo sharing can therefore be treated as relatively rich in everyday practice because it combines rapid interaction with visual content that can communicate relational and emotional information. At the same time, empirical work has qualified strong deterministic readings of MRT: Dennis and Kinney (1998) report that while cue and feedback properties shape perceptions of media richness, matching media richness to task equivocality did not reliably improve performance outcomes in their study \cite{dennis_kinney_1998}.

Social Presence Theory provides a complementary lens by emphasising the interpersonal qualities of mediated communication. Short, Williams, and Christie (1976) conceptualise social presence as the extent to which a medium conveys the salience of the interaction partner and the sense of interpersonal involvement, with higher social presence generally associated with greater warmth and immediacy in interaction \cite{short_williams_christie_1976}. In contemporary messaging contexts, this helps explain why users may prefer channels that support richer interpersonal signalling for relational communication, even when such channels also create additional privacy and control demands. Ishii et al. (2019) further argue that media richness should be revisited for modern technologies because users’ experiences of richness are shaped not only by channel properties but also by evolving features and patterns of use, which is particularly relevant for mobile messaging environments \cite{ishii_etal_2019}.

Taken together, MRT and Social Presence Theory help explain why revocable or dynamic consent in messaging is not only a technical problem but also a socially consequential one. When users select richer media for intimate or socially sensitive exchange, the interaction may carry stronger interpersonal involvement and relational expectations about how shared content ought to be handled \cite{daft_lengel_mrt_1986,short_williams_christie_1976}. In such contexts, withdrawing consent may be experienced as a delicate change in interpersonal boundaries rather than as a neutral technical action, because it can signal mistrust or create a sense of relational distance. Design implication: revocation mechanisms in WhatsApp-style photo sharing should be designed with attention to this social meaning, supporting face-saving and relationship maintenance while clearly communicating the practical limits of technical control in the presence of retention and recirculation practices \cite{dennis_kinney_1998,ishii_etal_2019}.

\section{The Concept of Consent}
This section examines consent in digital communication and argues that, in messaging contexts, it should be understood as an ongoing process rather than a one-time decision. While legal frameworks define the formal conditions under which consent is considered valid, consent in interpersonal messaging also functions as a relational and contextual expectation about what others may appropriately do with shared content.

\subsection{The Legal Mandate}
Although the GDPR is not designed specifically for consent revocation in online messaging platforms, it provides a useful benchmark for how valid consent and withdrawal should be implemented in such environments. Under the GDPR, consent must be freely given, specific, informed, and unambiguous, and withdrawal must be as easy as giving consent; where consent is withdrawn and no other legal basis remains, the right to erasure may also be engaged \cite{gdpr_2016,politou_2018}. Chaurasiya (2025) shows that the Indian DPDP Act adopts a similar consent-centred approach by defining consent in affirmative terms and recognising mechanisms for managing and withdrawing it \cite{chaurasiya_consent_2025}. However, as Politou et al. (2018) note, revocation operates prospectively rather than retroactively, and erasure becomes technically difficult once data have been copied, shared, or replicated across systems \cite{politou_2018}. In distributed communication contexts such as messaging, the legal mandate therefore establishes a right to revoke consent, but not a guarantee of complete practical forgetting.

\subsection{Consent Fatigue, the Validity Gap, and Non-Use}
Although legal and policy frameworks often treat consent as meaningful once it has been formally obtained, this assumption is difficult to sustain in contemporary digital environments. Baruh and Popescu argue that privacy self-management based on notice and choice places too much responsibility on individuals, even though data practices are often opaque, complex, and difficult to resist in practice \cite{baruh_big_2017}. In such settings, consent may be legally valid without being substantively meaningful, creating a validity gap between formal agreement and real user control. They further show that users may respond in two limited ways: through partial withdrawal or avoidance when privacy risks feel unmanageable, or through assimilation, where users accept reduced control and rely on market-provided protections \cite{baruh_big_2017}.

Waldman extends this critique by showing that privacy is often operationalised narrowly within organisations. In his account, lawyers tend to reduce privacy to notice and choice, while engineers often equate it with security or encryption, leaving limited space for more user-centred understandings of consent and control \cite{waldman_designing_nodate}. For the present topic, this suggests that effective revocation cannot depend only on formal consent or interface disclosure, but must be supported by designs that make control realistic and usable in practice.

\subsection{Consent Models}

Consent in everyday messaging cannot be understood adequately as a one-off, individualised act of permission, because interpersonal communication unfolds over time, across relationships, and under conditions of partial and often irreversible loss of control. Most of the revised consent models discussed below were developed in health, research, or broader data-governance settings rather than instant messaging; they are used here analytically, as conceptual tools for thinking about permission, revision, and withdrawal in WhatsApp-style communication. Traditional consent frameworks tend to assume that consent is static, singular, and transaction-like, yet this logic becomes unstable in ongoing interpersonal exchange, where the meaning of permission can shift after disclosure. As Burkhardt et al. suggest, meaningful consent depends not only on formal agreement but also on understanding and perceived control, while work on consent and revocation has also shown that decisions made at one moment cannot remain fully informed across the lifetime of data use \cite{burkhardt_privacy_2023,agrafiotis_consent_revocation_2012}. This is why dynamic consent has been positioned as a response to one-off authorisation models that ask people to agree to future uses they may neither anticipate nor revisit \cite{schuler_scott_dynamic_consent_2022}. In messaging, these weaknesses are particularly visible because permission is shaped by repeated interaction, changing context, recipient interpretation, and what happens after content has already been sent.

Dynamic Consent offers the clearest correction to this temporal problem by treating permission as revisable rather than fixed. Across the literature, it is characterised by ongoing engagement, granular choice, and the possibility of adjusting preferences over time rather than accepting a single all-or-nothing decision \cite{kaye_dynamic_2015,lee_identifying_2023}. This broader orientation is reinforced by Wilbanks’ argument that digital consent should be designed as the beginning of an ongoing ethical relationship rather than a one-time administrative step, and by later work that frames modifiable consent choices as central to user-centred data governance \cite{wilbanks_design_2018,narkhede_enhancing_2025}. However, what Dynamic Consent improves most clearly is temporality, not post-disclosure control. In everyday messaging, the ability to revise a preference does not necessarily translate into the ability to affect already-delivered content once it has moved beyond the sender’s immediate control.

Revocable Consent is therefore closer to the central concern of this dissertation, because it foregrounds withdrawal after permission has already been granted. The literature usefully distinguishes between revocability as a normative requirement and revocation as a practically enforceable outcome. Consent and revocation have been framed as ongoing controls, with informed revocation emphasising that users should understand the forms of withdrawal available to them \cite{agrafiotis_consent_revocation_2012}. Related technical work similarly treats withdrawal as a change in the scope of authorised access: users should be able to specify who may access data, for what purposes, and for how long, while formal models of consent evolution represent revocation by closing a previously open interval of collection or access rather than undoing prior disclosure \cite{merlec_smart_2021,robol_consent_2023}. This distinction is especially important for photo sharing, where a sender may be able to end future authorised access without being able to retract a photo that has already been viewed, saved, or copied by a recipient. Revocation in messaging is therefore best understood as limiting future legitimacy and platform-supported access, not as guaranteeing recovery of the original disclosure. This problem is reinforced by the broader notice-and-consent literature, which shows that even ostensibly reversible interfaces can fail to produce meaningful control when the design of choice itself is manipulative or unclear \cite{nouwens_dark_2020}.

Other models address narrower weaknesses in static consent without resolving the broader problem. Retroactive consent is useful because it makes the timing of authorisation explicit, allowing consent and withdrawal to be analysed in terms of whether their effects apply to past or future data practices \cite{robol_consent_2023}. A related response is to reduce persistence prospectively by limiting the duration or openness of permission. Although Im et al. do not formalise ephemeral consent as a standalone governance model, their emphasis on consent as revertible and specific supports the idea that permission should be bounded and revisable rather than indefinite \cite{im_yes_2021}. Yet these approaches remain partial. Once content has been copied, remembered, or redistributed by others, neither temporal expiry nor after-the-fact adjustment can fully restore sender control.

Seen comparatively, these models should not be treated as mutually exclusive competitors, but as attempts to repair different weaknesses in conventional consent. Dynamic Consent responds most clearly to temporality; Revocable Consent foregrounds withdrawal and reversibility; retroactive approaches sharpen how timing affects authorisation; and ephemeral approaches seek to limit persistence before problems arise. What unites them is a shared difficulty in accounting for messaging as a social practice in which the meaning and consequences of disclosure continue to evolve after the moment of sending \cite{kaye_dynamic_2015,robol_consent_2023,politou_forgetting_2018,im_yes_2021}. The central problem is therefore not simply lack of choice, but the challenge of aligning revisability, revocability, relationality, and downstream loss of control within a single model.

For this reason, affirmative and ongoing consent offers the strongest conceptual direction for this dissertation. Im et al. describe affirmative consent as voluntary, informed, revertible, specific, and unburdensome, providing a framework that treats permission as active and revisable rather than presumed or exhausted at a single point in time \cite{im_yes_2021}. These principles fit messaging more closely than one-off models because they better reflect how consent is negotiated in unfolding interaction. At the same time, they should not be overstated as a technical solution to screenshots, retention, or redistribution. Their value is conceptual and normative: they offer a clearer account of what responsible design should support, namely permission that can be revisited, clarified, limited, and renegotiated without assuming that post-sharing control can ever be complete.

This issue becomes even more complex in photo sharing, where consent may extend beyond a single sender-recipient relationship. 

Taken together, the literature suggests that the main weakness of conventional consent models is not simply insufficient granularity, but their inability to account for the ongoing social life of shared content. Improved models repair important dimensions of this problem, but they still encounter hard limits once messages and photos enter interpersonal networks where recipients can retain, reinterpret, copy, or redistribute them. For everyday messaging, meaningful consent is therefore best understood not as a completed transaction but as an ongoing and socially situated process whose practical force depends on both relational norms and the technical limits of post-sharing control.


\section{Consent in Instant Messaging and Photo Sharing}
This section connects the consent concepts and theories to the specific realities of photo sharing in instant messaging, focusing on WhatsApp context.

\subsection{Multi-Party Privacy Conflicts (MPCs)}
In this subsection, I will review literature on conflicts between uploaders/senders and co-owners/co-subjects of photos. The aim is to show that photo sharing frequently affects more than two parties: the sender, the recipient, and any depicted individuals. I will discuss Cherubini et al. (2021) to discuss the conflict between uploaders and data subjects and to compare precautionary mechanisms (prevent sharing until approved) versus dissuasion mechanisms (nudging or discouraging risky sharing). (In this chapter I use this distinction primarily to frame MPs; the HCI implications and interface-level operationalisation of these strategies are discussed in subsection 2.6.5) I will also use Xu et al. (2019) and Zhang et al. (2023) to discuss technical approaches for privacy-preserving photo sharing in social networks, using them to motivate why co-ownership cannot be solved purely by sender-side settings.

\subsection{Persistence, Recipient Control, and the Limits of Revocation}
In this subsection, I explain why consent revocation in instant messaging is structurally constrained by recipient-side control: messaging systems deliver content rather than enforce ongoing access control, so they cannot guarantee take-back after delivery.

Even when platforms provide deletion or ephemerality features, photos can remain accessible through device-level retention (e.g., local storage, notification previews, auto-download), re-sharing, or external capture (e.g., screenshots).
This creates a boundary condition for consent revocation: revocation cannot be treated as guaranteed erasure once a photo has been delivered to another person’s device.

I use WhatsApp as a concrete case to illustrate these constraints.
Although WhatsApp provides privacy-oriented features such as end-to-end encryption, photos are still delivered onto recipients’ devices and can therefore fall outside the sender’s control after sharing. As a result, revocation mechanisms in WhatsApp should be understood as limited post-delivery controls rather than guaranteed erasure, and interface design needs to communicate this limitation clearly to avoid misleading mental models (WhatsApp Help Center; WhatsApp Privacy Policy – UK). 
Deletion is the clearest example of this limitation: deletion delay means that “deleted” content can remain accessible (e.g., through lingering availability or indirect access) even after a sender attempts removal (Liang et al., 2015). WhatsApp states that recipients might see a message before it is deleted, and that deletion may not always be successful (WhatsApp Help Center), creating a mismatch between what users intend and what occurs (partial or failed removal).
Relatedly, WhatsApp provides disappearing messages and view-once media as privacy-oriented controls (WhatsApp Help Center), but these do not fully resolve persistence in practice because deletion and ephemerality can leave recoverable artefacts at the device/cloud level and can be bypassed through capture and re-sharing routes (Heath et al., 2023).

Deletion also has social trace effects: revocation becomes visible and socially meaningful when messaging interfaces signal removal (e.g., “This message was deleted”), which can create relational friction and boundary turbulence, especially in intimate contexts.

This raises a key question (“Even if revocation is possible, will people use it?”) because in interpersonal contexts, social pressure can discourage users from deleting or requesting deletion, particularly when deletion produces suspicion or conflict. 
Finally, I use Schnitzler et al. (2020) to show that users may misunderstand what deletion removes (local vs global deletion), producing overconfidence in control and inaccurate mental models of revocation outcomes (Schnitzler et al., 2020).
Therefore, revocation in WhatsApp photo sharing is an interface and expectation-management problem as much as a technical one: design must communicate what is and is not achievable, rather than implying that deletion produces universal forgetting.

\subsection{The Screenshot Threat}
In this subsection, I will review how screenshots and secondary sharing undermine revocation and ephemerality features, and how platform affordances shape both norms and misuse. While platforms introduce revocation-oriented controls (e.g., disappearing messages, view-once media), screenshots and external capture routes create a parallel persistence layer that exists outside of sender control and outside platform enforcement. WhatsApp’s view-once feature is explicitly designed to reduce retention and redistribution by preventing saving to the recipient’s gallery and limiting forwarding/copying within the app (WhatsApp Help Center). However, even where screenshot blocking is claimed or implemented in limited forms, capture and persistence can still occur through alternative pathways (e.g., other devices, screen recording in some contexts, or post-view replication), which means view-once operates more as a deterrent and norm signal than a guarantee of revocability.

To ground the screenshot issue socially (not only technically), but I will also use work on screenshot collection and sharing in messaging platforms to show that screenshots are often treated as a socially meaningful act that can violate privacy expectations and trigger CPM-style turbulence (Ingber, 2025). This supports the idea that screenshot risk is not merely a technical bypass but part of the interpersonal governance problem: users may assume contained sharing, while recipients may interpret screenshots as legitimate keeping, evidence, or sharing, producing conflict at the level of norms and expectations.

I will also connect screenshots and ephemeral messaging to evidence from forensic work showing that disappearing or ephemeral content can leave recoverable traces across devices and platforms, reinforcing that ephemerality does not necessarily eliminate persistence (Heath, 2023). In WhatsApp specifically, I will also note that the platform introduces partial safeguards such as chat-level restrictions (e.g., advanced chat privacy settings that limit certain forms of external sharing), but these remain limited relative to the broader ecosystem of copying, capture, and forwarding (WhatsApp Help Center). Overall, this subsection argues that the screenshot/secondary sharing threat is central to why consent revocation in WhatsApp photo sharing is difficult to operationalise as enforceable control, and why design must avoid overpromising revocation outcomes. 

\section{Human Computer Interaction (HCI) and Interface Design}
This section reviews HCI literature relevant to designing usable consent and revocation mechanisms, focusing on how design shapes understanding, decision-making, and control.
\subsection{Foundational Concepts and Goals of Usable Policy}
This subsection defines what usable privacy and consent controls mean in practice, focusing on how interface design shapes whether users can actually understand and exercise control. Building on Albesher and Alhussain (2021), usability in privacy contexts is treated as more than having settings available: controls must be findable, comprehensible, and actionable under everyday time and attention constraints, and they must communicate outcomes clearly enough to support accurate user mental models.

This subsection also draws on Privacy by Design to frame revocation as a core design responsibility rather than an afterthought. Cavoukian (2012) emphasises principles such as privacy as the default, proactive protection, and user-centricity, which are directly relevant to consent revocation because revocation is only meaningful when it is easy to access, simple to execute, and does not require expert knowledge. Taken together, these perspectives position consent revocation as a user-facing capability that must be designed for clarity, minimal friction, and realistic expectation-setting, rather than treated as a hidden compliance property that just exists “somewhere in settings.”

\subsection{Critical Failures of Current Interfaces}
This subsection examines the revocation gap, notice ineffectiveness, and design for obscurity (Albesher \& Alhussain, 2021; Bourgeus \& Vandercruysse, 2024; Schaub et al., 2018; Wilbanks, 2018). It focuses on how interface design can signal that control exists without supporting meaningful control in practice (e.g., vague labels, buried settings, unclear defaults, or ambiguous outcomes). This subsection also explains why mental models matter for revocation: if users misunderstand what a deletion or revocation feature actually changes, they cannot withdraw consent in an informed way.

\subsection{Manipulative Design}
This subsection examines dark patterns and manipulative consent interfaces (Bermejo Fernandez et al., 2021; Graßl et al.; Soe et al., 2020). It outlines common forms such as interface interference (Soe et al., 2020; Gray et al., 2021), obstruction and choice cascades (Soe et al., 2020), consent walls or forced action (Gray et al., 2021), unlabelled configuration sliders (Soe et al., 2020; Gray et al., 2021), and counterbalancing nudges (Bermejo Fernandez et al., 2021; Graßl et al.). The aim is to show how consent and control can be technically present while being practically discouraged, steered, or distorted through interaction design.

\subsection{Interactive Design Solutions for Meaningful Consent}
This subsection reviews interaction patterns intended to improve meaningful consent, focusing on how design can increase user understanding and reduce “automatic” agreement. It discusses the affirmative action requirement and the limitations of standard checkbox-based consent, drawing on Lindegren et al. (2021) to examine alternatives such as more engaged selection mechanisms (e.g., drag-and-drop or swipe actions) that can make the act of consenting more deliberate and potentially improve comprehension and intentionality. It also considers how the timing and format of consent interactions can affect whether users notice and interpret them during real tasks.

In addition, this subsection covers visual communication for clarity using Albesher and Alhussain (2021), highlighting how visual signifiers such as clearly differentiated states, consistent icons, and immediate feedback can improve users’ interpretation of privacy controls and reduce misunderstandings about what a setting will do. The overall focus is on interaction patterns that support informed choice without adding excessive friction or cognitive load.


\subsection{Designing for Revocation}
This subsection focuses on revocation-oriented design strategies by comparing precautionary mechanisms (e.g., pre-sharing approval or permission setting) with dissuasive mechanisms (e.g., nudges or warnings before sharing). It uses Schaub et al. (2018) to frame the design space for privacy notices particularly timing, channel, and modality and argues that just-in-time, contextual notices can support better decision-making than static policies, especially in photo-sharing moments where attention is limited and sharing is fast-paced.

\subsection{Operationalising Dynamic Consent in Messaging UI}
This subsection discusses dynamic consent as a design model that allows users to revise preferences over time, focusing on how ongoing control can be translated into messaging interfaces. Rather than repeating the conceptual definitions introduced in Section 2.4, it will concentrate on operationalising dynamic consent as usable, granular UI flows how permissions can be set, reviewed, updated, and understood in-context during everyday messaging.



\section{Empirical Usability Assessments of Current Platforms}
This section synthesises studies evaluating the usability of consent settings and revocation-related controls in WhatsApp and comparable platforms to ground the current state of practice. Albesher and Alhussain (2021) highlights severity usability problems affecting WhatsApp users’ ability to manage privacy controls, with particular emphasis on the STRAP heuristic (Specific, Transparent, Revocable, Actionable, Proportional) of the ability to revoke consent, which is reported as one of the most notable high-severity problem areas across the evaluated apps. Additionally, Kancherla et al. (2025)   discuss the technical reality of revocation and how platforms can fail to meet the practical expectations implied by revocability, including cases where revocation is harder than acceptance at the interface level and where revocation can fail “behind the UI” (e.g., not being properly stored, applied, or communicated). This supports the dissertation’s broader argument that “revocation availability” does not necessarily mean “revocation effectiveness.”

To connect these findings to WhatsApp’s concrete mechanisms, the section will review how revocation-oriented features currently operate and what they communicate to users. For example, WhatsApp’s deletion feature acknowledges that recipients may see a message before it is deleted and that deletion may not always be successful (WhatsApp, n.d.-a). It will also consider controls that aim to reduce secondary sharing, such as Advanced Chat Privacy, which is positioned as limiting content being taken outside WhatsApp (e.g., blocking chat export and disabling auto-download to gallery), while still operating within constraints where capture and manual retention can remain possible (WhatsApp, n.d.-b; WhatsApp, 2025). Finally, it will situate disappearing messages and view-once media as partial revocation-like mechanisms (WhatsApp, n.d.-c; WhatsApp, n.d.-d), while drawing on forensic and behavioural research suggesting that disappearing designs can still leave recoverable traces and that screenshot practices remain a persistent interpersonal privacy risk in messaging environments (Heath et al., 2023; Ingber, 2025). It will conclude by identifying where features such as delete for everyone, disappearing messages, and chat-level restrictions exist but do not constitute robust, granular, or reliably revertible consent mechanisms aligned with user expectations or the theoretical requirements discussed earlier.

\section{Research Gap and Summary}
This final section summarises unresolved issues and motivates the dissertation contribution.
\begin{itemize}
\item \textbf{Gap G1:} Interface-level translation of dynamic/ongoing consent remains underdeveloped in instant messaging.

Dynamic, ongoing, and revocable consent are discussed conceptually, but there is limited guidance on how to operationalise these models as usable UI flows within routine messaging.

\item \textbf{Gap G2:} Post-delivery controls in real platforms are partial, ambiguous, and prone to misleading mental models.

Features such as “delete for everyone” and disappearing options do not provide granular, policy-like control (e.g., per-photo expiry, download/forward restrictions) and can produce misleading mental models about what deletion achieves.

\item \textbf{Gap G3:} Multi-party privacy and co-subject agency are rarely supported in messenger UI.

Multi-party privacy conflicts show that depicted individuals may require agency, but most mainstream messenger tools assume only the sender controls the content.

\item \textbf{Gap G4:} Secondary sharing (forwarding/screenshots) undermines revocation, and designs often overpromise control.

Screenshot and forwarding practices make strict revocation difficult. Existing approaches provide partial deterrence rather than enforceable revocation.

\item \textbf{Gap G5:} Lack of agreed fair interaction models for sender vs recipient in revocation design.

Most work frames revocation primarily as a sender right, but less work specifies what “fair” and “intuitive” design means for both parties (e.g., notification, explanation, minimising blame while maintaining clarity).

\item \textbf{Gap G6:} Limited empirical understanding of user expectations for revocable photo permissions specifically in WhatsApp contexts.
Although privacy and safety concerns are documented, there is limited work on what WhatsApp users want from pre- and post-sharing control in everyday relational contexts, including the social friction and trace effects of revocation (e.g., deletion indicators).
\end{itemize}
This chapter therefore positions consent revocation as a socio-technical problem shaped by privacy theory, consent models, platform constraints, and interface design creating a clear need for user-centred revocation mechanisms for photo sharing in WhatsApp. 







